\documentclass[12pt]{amsart}
\usepackage[textwidth=13.5cm,textheight=21cm]{geometry}
\usepackage{amsmath,amssymb,amsthm,mathtools}
\usepackage[utf8]{inputenc}
\DeclareUnicodeCharacter{03B7}{\ensuremath{\eta}}
\DeclareUnicodeCharacter{03B5}{\ensuremath{\varepsilon}}
\DeclareUnicodeCharacter{211D}{\ensuremath{\mathbb{R}}}
\DeclareUnicodeCharacter{2115}{\ensuremath{\mathbb{N}}}
\DeclareUnicodeCharacter{2113}{\ensuremath{\ell}}
\DeclareUnicodeCharacter{2203}{\ensuremath{\exists}}
\DeclareUnicodeCharacter{2200}{\ensuremath{\forall}}
\DeclareUnicodeCharacter{2227}{\ensuremath{\wedge}}
\DeclareUnicodeCharacter{2264}{\ensuremath{\le}}
\DeclareUnicodeCharacter{2192}{\ensuremath{\to}}
\DeclareUnicodeCharacter{2080}{\ensuremath{_0}}
\DeclareUnicodeCharacter{03B8}{\ensuremath{\theta}}
\DeclareUnicodeCharacter{03C7}{\ensuremath{\chi}}
\DeclareUnicodeCharacter{2016}{\ensuremath{\Vert}}
\DeclareUnicodeCharacter{207B}{\ensuremath{^{-}}}
\DeclareUnicodeCharacter{00B9}{\ensuremath{^{1}}}
\DeclareUnicodeCharacter{208A}{\ensuremath{_{+}}}
\DeclareUnicodeCharacter{2102}{\ensuremath{\mathbb{C}}}
\DeclareUnicodeCharacter{2208}{\ensuremath{\in}}
\DeclareUnicodeCharacter{2211}{\ensuremath{\sum}}
\DeclareUnicodeCharacter{2260}{\ensuremath{\ne}}
\DeclareUnicodeCharacter{230A}{\ensuremath{\lfloor}}
\DeclareUnicodeCharacter{230B}{\ensuremath{\rfloor}}
\DeclareUnicodeCharacter{03B1}{\ensuremath{\alpha}}
\DeclareUnicodeCharacter{211A}{\ensuremath{\mathbb{Q}}}
\usepackage[hidelinks]{hyperref}
\usepackage{etoolbox}
\makeatletter
\patchcmd{\@setaddresses}{\footnotesize}{\footnotesize\raggedright}{}{}
\makeatother
\emergencystretch=3em
\newcommand{\zenodoID}{10.5281/zenodo.0000000}

\newtheorem{theorem}{Theorem}[section]
\newtheorem{lemma}[theorem]{Lemma}
\newtheorem{proposition}[theorem]{Proposition}
\newtheorem{corollary}[theorem]{Corollary}
\theoremstyle{remark}
\newtheorem{remark}[theorem]{Remark}

\renewcommand{\Re}{\operatorname{Re}}
\newcommand{\ord}{\operatorname{ord}}
\newcommand{\C}{\mathbb{C}}
\newcommand{\R}{\mathbb{R}}

\title{Vinogradov's conjectures on primitive roots and character sums}
\author{Deep Bhattacharjee}
\address{Formerly, Electro-Gravitational Space Propulsion Laboratory (EGSPL),
Bhubaneswar, Odisha 751030, India}
\email{d.bhattacharjee@erl-forschung.de, itsdeep@live.com}
\thanks{ORCID 0000-0003-0466-750X}
\subjclass[2020]{11A07, 11L40, 11M06, 11N13, 68V15}
\keywords{primitive roots, character sums, Dirichlet $L$-functions, zero-free regions, Lean 4}

\begin{document}

\begin{abstract}
For every $\eta>0$ every prime $p$ has a prime primitive root
$\ell\ll_\eta(\tau(p-1)^2\log p)^{8+\eta}$, so the least primitive root is $O_\varepsilon(p^\varepsilon)$.
For every nonprincipal character $\chi$ modulo $q$,
$\sum_{n\le x}\chi(n)\ll_\varepsilon q^\varepsilon x^{7/8+\varepsilon}$, which is $o(x)$ once
$x\ge q^\eta$. Both were conjectured by Vinogradov. For primitive $\chi$,
$\max_x|\sum_{n\le x}\chi(n)|\ll\sqrt q\log\log q$, which is optimal. The proofs use the zero-free
half-plane $\Re s>7/8$ and are checked in Lean~4, the last one with a theorem of Montgomery and
Vaughan as a hypothesis.
\end{abstract}

\maketitle

\section{Introduction}

For a prime $p$ let $g(p)$ be the least primitive root modulo $p$ and $g^*(p)$ the least prime
primitive root. Vinogradov conjectured that $g(p)\ll_\varepsilon p^\varepsilon$ for every
$\varepsilon>0$. Burgess~\cite{Bur62} proved $g(p)\ll_\varepsilon p^{1/4+\varepsilon}$; for the
prime version the least prime in a progression~\cite{HB92} gives $g^*(p)\ll p^{5.5}$. Under the
generalised Riemann hypothesis Shoup~\cite{Sho92} and Martin~\cite{Mar97} proved
$g^*(p)\ll\omega^4(\log\omega+1)^4\log^2p$ with $\omega=\omega(p-1)$. Erd\H{o}s~\cite{Erd65}
asked whether every prime $p$ has a prime $q<p$ that is a primitive root modulo $p$
(Problem~985 of~\cite{Blo}).

Let $\tau$ denote the divisor function. Throughout, $\chi$ is a Dirichlet character and
$L(s,\chi)$ its $L$-function.

\begin{theorem}\label{thm:main}
For every $\eta>0$ there is $C_\eta$ such that every prime $p$ has a prime primitive root $\ell$
with
\[
  \ell\le C_\eta\bigl(\tau(p-1)^2\log(p+2)\bigr)^{8+\eta}.
\]
\end{theorem}

\begin{corollary}\label{cor:vinogradov}
For every $\varepsilon>0$ there is $C_\varepsilon$ such that every prime $p$ has a prime primitive
root $\ell\le C_\varepsilon p^\varepsilon$. In particular $g(p)\le g^*(p)\ll_\varepsilon p^\varepsilon$.
\end{corollary}

\begin{corollary}\label{cor:erdos}
There is $p_0$ such that every prime $p\ge p_0$ has a prime $q<p$ that is a primitive root
modulo $p$.
\end{corollary}

Vinogradov also conjectured that $\sum_{n\le x}\chi(n)=o(x)$ for nonprincipal characters
$\chi$ modulo $q$ once $x\ge q^\varepsilon$; see~\cite{GS18}. For prime moduli Burgess~\cite{Bur62}
proved this for $x\ge q^{1/4+\varepsilon}$, and it follows from the generalised Riemann hypothesis.

\begin{theorem}\label{thm:lindelof}
For every $\delta>0$ there is $C_\delta$ such that for every nonprincipal character $\chi$ modulo
$q$ and every $s=\sigma+it$ with $\sigma\ge7/8+\delta$,
\[
  |L(s,\chi)|+|L(s,\chi)|^{-1}\le C_\delta\bigl(q(3+|t|)\bigr)^\delta.
\]
\end{theorem}

\begin{theorem}\label{thm:charsum}
For every $\delta>0$ there is $C_\delta$ such that for every nonprincipal character $\chi$ modulo
$q$ and every $x\ge1$,
\[
  \Bigl|\sum_{n\le x}\chi(n)\Bigr|\le C_\delta\,q^\delta x^{7/8+\delta}.
\]
\end{theorem}

\begin{corollary}\label{cor:charsum}
For every $\eta>0$ and $\theta>7/8$ there is $C$ such that $\bigl|\sum_{n\le x}\chi(n)\bigr|\le
Cx^\theta$ for every nonprincipal character $\chi$ modulo $q$ and every $x\ge q^\eta$. In
particular $\sum_{n\le x}\chi(n)=o(x)$ uniformly for $x\ge q^\eta$.
\end{corollary}

The P\'olya--Vinogradov inequality~\cite[Ch.~23]{Dav00} gives
$\bigl|\sum_{n\le x}\chi(n)\bigr|\ll\sqrt q\log q$ for every nonprincipal $\chi$ modulo $q$. Paley~\cite{Pal32}
showed that $\log\log q$ cannot replace $\log q$ with a factor tending to $0$: for infinitely many
primitive quadratic characters the maximum exceeds $c\sqrt q\log\log q$. Montgomery and
Vaughan~\cite{MV77} proved the upper bound $\sqrt q\log\log q$ for primitive characters under the
generalised Riemann hypothesis, and Granville and Soundararajan~\cite{GS07} improved the
exponent of $\log q$ unconditionally for characters of fixed odd order.

\begin{theorem}\label{thm:pv}
There is an absolute constant $C$ such that for every primitive character $\chi$ modulo $q\ge3$
and every $x\ge1$,
\[
  \Bigl|\sum_{n\le x}\chi(n)\Bigr|\le C\sqrt q\log\log q.
\]
\end{theorem}

Montgomery and Vaughan use the generalised Riemann hypothesis only for prime sums of the
characters $\chi\psi$, with $\psi$ modulo a small $r$, on the major arcs; Theorem~\ref{thm:zerofree}
below suffices there. On the minor arcs they use their unconditional bound for exponential sums
with multiplicative coefficients, Theorem~\ref{thm:mv}.

The input is the zero-free half-plane of~\cite{OAI26}, from which~\cite{OAI26} deduces, as
in~\cite[Thm.~6.7]{BIMS17}, that the least quadratic nonresidue is $\ll_\varepsilon p^\varepsilon$.
A primitive root must avoid every subgroup of prime index at once, and the number of characters
involved is the source of the factor $\tau(p-1)^2$.

\begin{theorem}[{\cite{OAI26}}]\label{thm:zerofree}
If $\chi$ is a nonprincipal Dirichlet character, then $L(s,\chi)\ne0$ for $\Re s>7/8$.
\end{theorem}

Theorem~\ref{thm:zerofree} is formalised in Lean~4 in~\cite{OAI26}. Theorems~\ref{thm:main},
\ref{thm:lindelof} and~\ref{thm:charsum} and Corollaries~\ref{cor:vinogradov}, \ref{cor:erdos}
and~\ref{cor:charsum} are formalised here, with Theorem~\ref{thm:zerofree} as the only external
input, and Theorem~\ref{thm:pv} is formalised from Theorems~\ref{thm:zerofree} and~\ref{thm:mv}
(Section~\ref{sec:lean}). The constants are effective.

Theorem~\ref{thm:lindelof} is Littlewood's deduction of a Lindel\"of bound from a zero-free
half-plane~\cite[\S14.2]{Tit86}, and Theorem~\ref{thm:charsum} follows from it by a contour shift;
the new ingredient is Theorem~\ref{thm:zerofree}. Among the consequences of
Theorem~\ref{thm:zerofree}, \cite{OAI26} lists the least quadratic nonresidue, the prime number
theorem in progressions, class numbers and idoneal numbers, but not
Theorems~\ref{thm:lindelof}, \ref{thm:charsum} and~\ref{thm:pv}. All results here hold with $7/8$ replaced by any
$\theta\in[1/2,1)$ such that $L(s,\chi)\ne0$ for $\Re s>\theta$ and every nonprincipal $\chi$, for
instance by the value $0.874957\ldots$ claimed in~\cite{Liu26}. The companion
papers~\cite{OAI26b,OAI26c} prove, from a zero-free strip $\Re s>1-10^{-6}$ for Hecke
$L$-functions of cyclotomic fields containing the twelfth roots of unity, that every integer other
than $-1$ and the squares is a primitive root modulo infinitely many primes, and that polynomials
over $\mathbb F_p$ can be factored in deterministic polynomial time; neither bounds $g(p)$.

\begin{remark}\label{rem:p0}
The constant $p_0$ of Corollary~\ref{cor:erdos} is effective but far beyond direct computation.
A computation (Section~\ref{sec:lean}) gives $g^*(p)<p$ for every prime $3\le p\le10^8$, with
$\max g^*(p)=223$ at $p=73189117$. Erd\H{o}s's question therefore remains open for the primes in
$[10^8,p_0)$.
\end{remark}

\begin{remark}\label{rem:wigert}
Since $\tau(n)\le 2^{(1+o(1))\log n/\log\log n}$ \cite[Thm.~317]{HW}, Theorem~\ref{thm:main}
gives $g^*(p)\le\exp\bigl(c\log p/\log\log p\bigr)$ for an absolute constant $c$. This form is not
formalised.
\end{remark}

\begin{remark}\label{rem:carella}
Carella has claimed $g^*(p)\ll p^{5/\log\log p}$ \cite{Car15} and
$g^*(p)=O((\log p)^{1+\varepsilon})$ \cite{Car18,Car25} without hypotheses. In~\cite{Car15}, the
estimate~(16) bounds $\sum_{n\le x}\chi(n)\Lambda(n)n^{-2}$ by $O(1/x)$; the sum tends to
$-L'(2,\chi)/L(2,\chi)$, which equals $-0.17264\ldots$ for the character modulo~$3$. In the
arXiv version of~\cite{Car18}, Lemma~7.1 asserts $E(x,y)\ll yp^{-\varepsilon}$ for the sum
$E(x,y)$ of~(7.2) when no integer of $[x,x+y]$ is a primitive root. Under that hypothesis every
inner sum over $m$ equals $-1$, so $E(x,y)=-\varphi(p-1)\,\#([x,x+y]\cap\mathbb Z)/p$ exactly,
while $\varphi(p-1)/p\gg1/\log\log p$; the step~(7.2) replaces $\sum_ne(m\tau^n/p)$ by its value
at $m=1$. We could not obtain the text of~\cite{Car25}.
\end{remark}

\section{The prime sum of a nonprincipal character}

For an arithmetic function $f$ and $x\ge1$ put
\[
  S(f;x)=\sum_{n\le x}f(n)\Lambda(n)\Bigl(1-\frac nx\Bigr),
\]
and write $S(\chi;x)$ when $f=\chi$. The aim of this section is Proposition~\ref{prop:riesz}.

\begin{lemma}\label{lem:mellin}
Let $a\colon\mathbb N\to\C$ satisfy $\sum_n|a(n)|n^{-2}<\infty$ and put $F(s)=\sum_na(n)n^{-s}$
for $\Re s\ge2$. For $x>0$,
\[
  \sum_{n\le x}a(n)\Bigl(1-\frac nx\Bigr)=\frac1{2\pi}\int_{-\infty}^{\infty}
  F(2+it)\,\frac{x^{2+it}}{(2+it)(3+it)}\,dt.
\]
\end{lemma}

\begin{proof}
For $y>0$ let $k(y)=\frac1{2\pi i}\int_{(2)}y^s\,\frac{ds}{s(s+1)}$. If $y\ge1$, moving the line
to the left through the poles at $0$ and $-1$ gives $k(y)=1-1/y$; if $y<1$, moving it to the right
gives $k(y)=0$. Since $|s(s+1)|^{-1}$ is integrable on $\Re s=2$ and
$\sum_n|a(n)|n^{-2}<\infty$, summation and integration may be exchanged:
the right side equals $\sum_na(n)k(x/n)$.
\end{proof}

\begin{lemma}\label{lem:size}
Let $\chi$ be nonprincipal modulo $q$. For $\Re s=\sigma>0$, $|L(s,\chi)|\le q|s|/\sigma$.
\end{lemma}

\begin{proof}
Put $A(t)=\sum_{n\le t}\chi(n)$; then $|A(t)|\le q$, and partial summation gives
$L(s,\chi)=s\int_1^\infty A(t)t^{-s-1}\,dt$ for $\sigma>0$. Hence
$|L(s,\chi)|\le q|s|\int_1^\infty t^{-\sigma-1}dt=q|s|/\sigma$.
\end{proof}

\begin{lemma}\label{lem:euler}
Let $A_0=\frac32\sum_{\ell}\ell^{-3/2}$, the sum over primes. For every Dirichlet character
$\chi$ the series $E(s)=\sum_\ell-\log(1-\chi(\ell)\ell^{-s})$ is holomorphic on $\Re s>1$,
$e^{E(s)}=L(s,\chi)$ there, and $|E(s)|\le A_0$ for $\Re s\ge3/2$.
\end{lemma}

\begin{proof}
For $\Re s\ge1$ we have $|\chi(\ell)\ell^{-s}|\le1/2$, and $|\log(1-w)|\le\frac32|w|$ for
$|w|\le1/2$. The series therefore converges locally uniformly on $\Re s>1$ and is bounded by
$A_0$ on $\Re s\ge3/2$, and the Euler product gives $e^{E}=L$.
\end{proof}

\begin{lemma}\label{lem:logderiv}
Let $0<\delta<1/1000$, $1/2\le a\le1$, and let $\chi$ be nonprincipal modulo $q$ with
$L(s,\chi)\ne0$ for $\Re s>a+2\delta$. Put $K_1=(1+2A_0)/(2\delta^2)$. For
$a+8\delta\le\sigma\le2$ and $t\in\R$,
\[
  \Bigl|\frac{L'}{L}(\sigma+it,\chi)\Bigr|\le K_1\log\bigl(2q(3+|t|)^2\bigr).
\]
\end{lemma}

\begin{proof}
Let $c=2+it$, $R=2-a-3\delta$, $r=2-a-6\delta$ and $\Lambda_0=\log(2q(3+|t|)^2)\ge\log18>1$.
On the disc $|z-c|<2-a-2\delta$ we have $\Re z>a+2\delta$, so $L(z,\chi)$ has no zero there
and has a holomorphic logarithm $h$. On $|z-c|<1/2$ both $h$ and $E$ of Lemma~\ref{lem:euler}
are logarithms of $L$, so they differ by a constant in $2\pi i\mathbb Z$; we choose $h=E$ there.
Then $|h(c)|\le A_0$.

If $|z-c|\le R$, then $\Re z\ge a+3\delta\ge1/2$ and $|z|\le|c|+3/2\le4+|t|$, so
Lemma~\ref{lem:size} gives $\Re h(z)=\log|L(z,\chi)|\le\log(2q(4+|t|))\le\Lambda_0$.
The Borel--Carath\'eodory inequality~\cite[\S5.5]{Tit39} gives, for $|z-c|\le r$,
\[
  |h(z)-h(c)|\le\frac{2r}{R-r}\max_{|w-c|=R}\Re\bigl(h(w)-h(c)\bigr)
  \le\frac{\Lambda_0+A_0}{\delta},
\]
as $r\le3/2$ and $R-r=3\delta$. Hence $|h(z)|\le A_0+(\Lambda_0+A_0)/\delta\le(1+2A_0)\Lambda_0/\delta$
on $|z-c|\le r$. For $w=\sigma+it$ with $a+8\delta\le\sigma\le2$ we have $|w-c|=2-\sigma\le r-2\delta$,
so the circle $|z-w|=2\delta$ lies in $|z-c|\le r$, and Cauchy's estimate gives
$|L'/L(w,\chi)|=|h'(w)|\le(1+2A_0)\Lambda_0/(2\delta^2)$.
\end{proof}

\begin{lemma}\label{lem:elementary}
For $q\ge1$, $t\in\R$ and $s=\sigma+it$ with $\sigma\ge1/2$,
\[
  \log\bigl(2q(3+|t|)^2\bigr)\le10\log(q+2)(1+|t|)^{1/2},\qquad |s(s+1)|\ge\tfrac18(1+|t|)^2.
\]
\end{lemma}

\begin{proof}
We have $\log(2q)\le2\log(q+2)$, and $\log y\le2y^{1/2}$ with $3+|t|\le4(1+|t|)$ gives
$\log(3+|t|)\le4(1+|t|)^{1/2}$; the first bound follows from $\log(q+2)\ge1$. For the second,
$|s|\ge\max(\sigma,|t|)\ge(1+|t|)/4$ and $|s+1|\ge\max(\sigma+1,|t|)\ge(1+|t|)/2$.
\end{proof}

\begin{proposition}\label{prop:riesz}
Let $0<\delta<1/1000$. There is $K=K(\delta)$ such that the following holds. Let $1/2\le a\le1$
and let $\chi$ be nonprincipal modulo $q$ with $L(s,\chi)\ne0$ for $\Re s>a+2\delta$. Then for
$x\ge1$,
\[
  |S(\chi;x)|\le K\,x^{a+8\delta}\log(q+2).
\]
\end{proposition}

\begin{proof}
Put $\sigma_0=a+8\delta$ and
\[
  \Phi(s)=-\frac{L'}{L}(s,\chi)\,\frac{x^s}{s(s+1)}.
\]
For $\Re s>1$, $-L'/L(s,\chi)=\sum_n\chi(n)\Lambda(n)n^{-s}$, so Lemma~\ref{lem:mellin} gives
$S(\chi;x)=\frac1{2\pi}\int\Phi(2+it)\,dt$. The function $\Phi$ is holomorphic on
$\sigma_0\le\Re s\le2$, and by Lemmas~\ref{lem:logderiv} and~\ref{lem:elementary},
\[
  |\Phi(\sigma+it)|\le80K_1\log(q+2)\,x^{\sigma}(1+|t|)^{-3/2}\qquad(\sigma_0\le\sigma\le2).
\]
This tends to $0$ as $|t|\to\infty$ uniformly in $\sigma$, so Cauchy's theorem on the rectangles
$[\sigma_0,2]\times[-T,T]$, $T\to\infty$, moves the integral to $\Re s=\sigma_0$. Since
$\int_\R(1+|t|)^{-3/2}dt=4$, we get $|S(\chi;x)|\le\frac{160K_1}{\pi}x^{\sigma_0}\log(q+2)$.
\end{proof}

\section{Primitive roots as a combination of characters}

Fix a prime $p$, put $N=p-1$, let $g$ generate $(\mathbb Z/p\mathbb Z)^\times$, and let $\chi_1$
be the character modulo $p$ with $\chi_1(g)=e^{2\pi i/N}$. Let $\chi_0$ be the principal character.

\begin{lemma}\label{lem:indicator}
For every integer $n$,
\[
  \mathbf 1[n\text{ is a primitive root mod }p]=\sum_{d\mid N}\frac{\mu(d)}d\sum_{m=0}^{d-1}
  \chi_1^{mN/d}(n),
\]
and $\chi_1^{mN/d}\ne\chi_0$ for $d\mid N$ and $0<m<d$.
\end{lemma}

\begin{proof}
If $p\mid n$ both sides vanish. Otherwise $n\equiv g^k$, the order of $n$ is $N/(N,k)$, and
$\sum_{m<d}e^{2\pi ikm/d}=d\,\mathbf 1[d\mid k]$; the right side is
$\sum_{d\mid(N,k)}\mu(d)=\mathbf 1[(N,k)=1]$. The character $\chi_1^{mN/d}$ sends $g$ to
$e^{2\pi im/d}\ne1$.
\end{proof}

\begin{lemma}\label{lem:arith}
For $N\ge1$, $\sum_{d\mid N}\mu(d)/d=\varphi(N)/N\ge1/\tau(N)$.
\end{lemma}

\begin{proof}
$\varphi(N)/N=\prod_{r\mid N}(1-1/r)\ge2^{-\omega(N)}\ge1/\tau(N)$, the product over the prime
divisors $r$ of $N$.
\end{proof}

\begin{lemma}\label{lem:principal}
For $x\ge10^4$, $\Re S(\chi_0;x)\ge\frac12\psi(x/2)-\log x\ge x/10$.
\end{lemma}

\begin{proof}
The terms of $\sum_{n\le x}\Lambda(n)(1-n/x)$ are nonnegative and at least $\Lambda(n)/2$ for
$n\le x/2$, so the sum is at least $\psi(x/2)/2$. The terms removed by $\chi_0$ are those with
$p\mid n$, that is $n=p^k\le x$, and they contribute at most $\sum_{p^k\le x}\log p\le\log x$.
For the second inequality, $\psi(n)\ge n\log2$ for integers $n\ge7$~\cite{Nai82}, so
$\psi(x/2)\ge(x/2-1)\log2$, and $\frac x4\log2-\frac12\log2-\log x\ge\frac x{10}$ for $x\ge10^4$.
\end{proof}

\begin{lemma}\label{lem:powers}
If $\ell^k$ with $k\ge1$ is a primitive root modulo $p$, so is $\ell$.
\end{lemma}

\begin{proof}
$\ell$ is a unit modulo $p$, and $N=\ord(\ell^k)$ divides $\ord(\ell)$, which divides $N$.
\end{proof}

\section{Proof of Theorem~\ref{thm:main} and Corollaries~\ref{cor:vinogradov} and~\ref{cor:erdos}}

\begin{proof}[Proof of Theorem~\ref{thm:main}]
Let $\delta=\eta/(2000(8+\eta))$, so $0<\delta<1/1000$, and let $K$ be given by
Proposition~\ref{prop:riesz}; put $K'=\max(K,1)$. Let $\sigma_0=7/8+8\delta$ and
$\beta=1-\sigma_0=1/8-8\delta>0$. Then
\[
  (8+\eta)\beta=1+\frac\eta8-\frac\eta{250}\ge1.
\]
Put $C_\eta=\max\bigl(10^4,(20K')^{8+\eta}\bigr)$. Given a prime $p$, let $\tau=\tau(p-1)$,
$Y=\tau^2\log(p+2)\ge1$ and $x=C_\eta Y^{8+\eta}\ge10^4$.

By Theorem~\ref{thm:zerofree}, Proposition~\ref{prop:riesz} applies with $a=7/8$ to every
nonprincipal $\chi$ modulo $p$, so $|S(\chi;x)|\le B:=K'x^{\sigma_0}\log(p+2)$. Let $G$ be the sum
$S(f;x)$ with $f$ the indicator of the primitive roots modulo $p$. By Lemma~\ref{lem:indicator},
\[
  G=c_0S(\chi_0;x)+\mathcal E,\qquad c_0=\sum_{d\mid N}\frac{\mu(d)}d,\qquad
  \mathcal E=\sum_{d\mid N}\frac{\mu(d)}d\sum_{0<m<d}S(\chi_1^{mN/d};x),
\]
and $|\mathcal E|\le\sum_{d\mid N}\frac{d-1}dB\le\tau B$. Lemmas~\ref{lem:arith}
and~\ref{lem:principal} give
\[
  |G|\ge c_0\Re S(\chi_0;x)-|\mathcal E|\ge\frac x{10\tau}-\tau B.
\]
Now $x^\beta=C_\eta^\beta Y^{(8+\eta)\beta}\ge(20K')^{(8+\eta)\beta}Y\ge20K'Y$, so
$\tau^2B=K'Yx^{\sigma_0}\le x^{\beta+\sigma_0}/20=x/20$. Hence $|G|\ge x/(20\tau)>0$, and some
$n\le x$ with $\Lambda(n)\ne0$ is a primitive root modulo $p$. Writing $n=\ell^k$, the prime
$\ell\le x$ is a primitive root by Lemma~\ref{lem:powers}.
\end{proof}

\begin{proof}[Proof of Corollary~\ref{cor:vinogradov}]
Take $\eta=1$ and $\nu=\varepsilon/27$. By~\cite[Thm.~315]{HW} there is $D_\nu$ with
$\tau(n)\le D_\nu n^\nu$, and $\log(p+2)\le(p+2)^\nu/\nu\le3^\nu p^\nu/\nu$. Theorem~\ref{thm:main}
gives $\ell\le C_1(D_\nu^2 3^\nu\nu^{-1})^9p^{27\nu}$.
\end{proof}

\begin{proof}[Proof of Corollary~\ref{cor:erdos}]
Corollary~\ref{cor:vinogradov} with $\varepsilon=1/2$ gives $\ell\le C\sqrt p<p$ for
$p>C^2$.
\end{proof}

\section{The Lindel\"of bound and character sums}\label{sec:charsum}

In this section $A_0$ is the constant of Lemma~\ref{lem:euler}.

\begin{lemma}\label{lem:lindelof}
Let $0<\kappa<1/1000$ and $\varepsilon>0$. There is $D=D(\kappa,\varepsilon)$ such that the
following holds. Let $1/2\le a\le1$ and let $\chi$ be nonprincipal modulo $q$ with
$L(s,\chi)\ne0$ for $\Re s>a+2\kappa$. For $a+8\kappa\le\sigma\le2$ and $t\in\R$,
\[
  |L(\sigma+it,\chi)|+|L(\sigma+it,\chi)|^{-1}\le D\bigl(2q(3+|t|)^2\bigr)^\varepsilon.
\]
\end{lemma}

\begin{proof}
Let $c$, $r$, $\Lambda_0$ and the logarithm $h$ of $L(z,\chi)$ on $|z-c|<2-a-2\kappa$ be as in
the proof of Lemma~\ref{lem:logderiv}, with $\delta$ replaced by $\kappa$. Both $h$ and $E$ are
logarithms of $L$ on the disc $|z-c|<3/4$, where $\Re z>1$, so $h=E$ there, and
$|h(z)|\le A_0$ for $|z-c|\le1/2$ by Lemma~\ref{lem:euler}. The proof of
Lemma~\ref{lem:logderiv} gives $|h(z)|\le M:=(1+2A_0)\Lambda_0/\kappa$ for $|z-c|\le r$.

Let $w=\sigma+it$, $\rho=|w-c|=2-\sigma$, $A_1=\max(A_0,1)$ and $\lambda=4\kappa/(3\log3)$; note
$M\ge1$. If $\rho\le1/2$ then $|h(w)|\le A_0\le A_1M^{1-\lambda}$. Otherwise
$1/2<\rho\le r-2\kappa$, and Hadamard's three circles theorem~\cite[\S5.3]{Tit39} gives
$|h(w)|\le A_0^{1-\alpha}M^{\alpha}$ with $\alpha=\log(2\rho)/\log(2r)$. Since
$1-6\kappa\le r\le3/2$,
\[
  \alpha\le\frac{\log(2r-4\kappa)}{\log(2r)}=1+\frac{\log(1-2\kappa/r)}{\log(2r)}
  \le1-\frac{2\kappa/r}{\log3}\le1-\lambda,
\]
so again $|h(w)|\le A_1M^{1-\lambda}=B\Lambda_0^{1-\lambda}$ with
$B=A_1((1+2A_0)/\kappa)^{1-\lambda}$. If $\Lambda_0^\lambda\ge B/\varepsilon$ then
$B\Lambda_0^{1-\lambda}\le\varepsilon\Lambda_0$, and otherwise $B\Lambda_0^{1-\lambda}\le
D_1:=B(B/\varepsilon)^{(1-\lambda)/\lambda}$. Since $\log|L(w,\chi)|=\Re h(w)$,
\[
  |L(w,\chi)|^{\pm1}\le e^{|h(w)|}\le e^{D_1}e^{\varepsilon\Lambda_0}
  =e^{D_1}\bigl(2q(3+|t|)^2\bigr)^\varepsilon,
\]
and $D=2e^{D_1}$ works.
\end{proof}

\begin{lemma}\label{lem:unsmooth}
Let $|b(n)|\le1$ for $n\ge1$ and put $T(y)=\sum_{n\le y}b(n)(y-n)$ for $y>0$. For $x,h>0$,
\[
  h\Bigl|\sum_{n\le x}b(n)\Bigr|\le|T(x+h)-T(x)|+h(h+1).
\]
\end{lemma}

\begin{proof}
$T(x+h)-T(x)=h\sum_{n\le x}b(n)+\sum_{x<n\le x+h}b(n)(x+h-n)$, and the last sum has at most
$h+1$ terms, each of modulus at most $h$.
\end{proof}

\begin{proposition}\label{prop:charsum}
Let $0<\kappa<1/1000$ and $0<\varepsilon\le1/8$. There is $K=K(\kappa,\varepsilon)$ such that the
following holds. Let $1/2\le a$ with $a+8\kappa\le1$, and let $\chi$ be nonprincipal modulo $q$
with $L(s,\chi)\ne0$ for $\Re s>a+2\kappa$. Then for $x\ge1$,
\[
  \Bigl|\sum_{n\le x}\chi(n)\Bigr|\le K\,q^\varepsilon x^{a+8\kappa+4\varepsilon}.
\]
\end{proposition}

\begin{proof}
Put $\sigma_0=a+8\kappa\in[1/2,1]$ and let $D$ be given by Lemma~\ref{lem:lindelof}. Since
$2q(3+|t|)^2\le18q(1+|t|)^2$, Lemmas~\ref{lem:lindelof} and~\ref{lem:elementary} give, for
$s=\sigma+it$ with $\sigma_0\le\sigma\le2$,
\begin{equation}\label{eq:quot}
  \Bigl|\frac{L(s,\chi)}{s(s+1)}\Bigr|\le8D(18q)^\varepsilon(1+|t|)^{-2+2\varepsilon}.
\end{equation}
Let $T(y)=\sum_{n\le y}\chi(n)(y-n)$. Lemma~\ref{lem:mellin} with $a(n)=\chi(n)$ gives
$T(y)=\frac1{2\pi i}\int_{(2)}L(s,\chi)\,y^{s+1}\frac{ds}{s(s+1)}$. The integrand is entire, and
by~\eqref{eq:quot} it tends to $0$ as $|t|\to\infty$ uniformly in $\sigma_0\le\sigma\le2$, since
$2\varepsilon<1$; Cauchy's theorem moves the line to $\Re s=\sigma_0$. For $h>0$,
\[
  T(x+h)-T(x)=\frac1{2\pi}\int_{-\infty}^{\infty}\frac{L(s,\chi)}{s(s+1)}
  \bigl((x+h)^{s+1}-x^{s+1}\bigr)\,dt,\qquad s=\sigma_0+it.
\]
Since $(x+h)^{s+1}-x^{s+1}=(s+1)\int_x^{x+h}u^s\,du$ and $|s+1|\le3(1+|t|)$, its modulus is at
most $3(1+|t|)h(x+h)^{\sigma_0}$, and also at most $2(x+h)^{\sigma_0+1}$. The minimum of two
numbers is at most the product of their powers $1-4\varepsilon$ and $4\varepsilon$, so
\[
  \bigl|(x+h)^{s+1}-x^{s+1}\bigr|\le(1+|t|)^{1-4\varepsilon}(x+h)^{\sigma_0}(3h)^{1-4\varepsilon}
  \bigl(2(x+h)\bigr)^{4\varepsilon}.
\]
With~\eqref{eq:quot} and $\int_\R(1+|t|)^{-1-2\varepsilon}dt=1/\varepsilon$,
\[
  |T(x+h)-T(x)|\le\frac{4D(18q)^\varepsilon}{\pi\varepsilon}(x+h)^{\sigma_0}(3h)^{1-4\varepsilon}
  \bigl(2(x+h)\bigr)^{4\varepsilon}.
\]
Take $h=x^{\sigma_0}$, so $1\le h\le x$. Then $(x+h)^{\sigma_0}\le2h$, $(3h)^{1-4\varepsilon}\le3h$
and $(2(x+h))^{4\varepsilon}\le4x^{4\varepsilon}$, so the right side is at most
$\frac{1728D}{\pi\varepsilon}q^\varepsilon h^2x^{4\varepsilon}$. By Lemma~\ref{lem:unsmooth}
with $b=\chi$,
\[
  \Bigl|\sum_{n\le x}\chi(n)\Bigr|\le\frac{1728D}{\pi\varepsilon}q^\varepsilon hx^{4\varepsilon}+h+1
  \le\Bigl(\frac{1728D}{\pi\varepsilon}+2\Bigr)q^\varepsilon x^{\sigma_0+4\varepsilon}.
  \qedhere
\]
\end{proof}

\begin{proof}[Proof of Theorem~\ref{thm:lindelof}]
Let $\kappa=\min(\delta/8,1/2000)$ and $a=7/8-2\kappa$, so $1/2\le a\le1$ and, by
Theorem~\ref{thm:zerofree}, $L(s,\chi)\ne0$ for $\Re s>a+2\kappa$. Let $D$ be given by
Lemma~\ref{lem:lindelof} with $\varepsilon=\delta/2$. If $7/8+\delta\le\sigma\le2$, then
$\sigma\ge a+8\kappa$, and $2q(3+|t|)^2\le(2q(3+|t|))^2$ gives
$|L|+|L|^{-1}\le D2^\delta(q(3+|t|))^\delta$. If $\sigma>2$, Lemma~\ref{lem:euler} gives
$|L|^{\pm1}\le e^{A_0}$. So $C_\delta=D2^\delta+2e^{A_0}$ works.
\end{proof}

\begin{proof}[Proof of Theorem~\ref{thm:charsum}]
Let $\delta'=\min(\delta,1)$, $\kappa=\min(\delta'/12,1/2000)$, $a=7/8-2\kappa$ and
$\varepsilon=\delta'/8$. Then $a+8\kappa=7/8+6\kappa\le1$, and Proposition~\ref{prop:charsum},
with Theorem~\ref{thm:zerofree}, gives the bound
$Kq^{\delta'/8}x^{7/8+6\kappa+\delta'/2}\le Kq^\delta x^{7/8+\delta}$.
\end{proof}

\begin{proof}[Proof of Corollary~\ref{cor:charsum}]
Let $\delta=(\theta-7/8)\eta/(1+\eta)$, so that $\delta/\eta+7/8+\delta=\theta$. For
$x\ge q^\eta$ we have $q^\delta\le x^{\delta/\eta}$, and Theorem~\ref{thm:charsum} gives
$|\sum_{n\le x}\chi(n)|\le C_\delta x^\theta$. With $\theta=15/16$ this is $o(x)$.
\end{proof}


\section{The P\'olya--Vinogradov inequality}\label{sec:pv}

Write $e(x)=e^{2\pi ix}$, and let $P(n)$ be the largest prime factor of $n\ge2$, with $P(1)=1$.
The minor arcs are handled by the following theorem.

\begin{theorem}[{\cite[Corollary~1]{MV77}}]\label{thm:mv}
There is an absolute constant $C_1$ with the following property. Let $f$ be multiplicative with
$|f(n)|\le1$, let $\alpha\in\R$, and let $a,r$ be coprime integers with
$|\alpha-a/r|\le r^{-2}$ and $2\le R\le r\le N/R$. Then
\[
  \Bigl|\sum_{n\le N}f(n)e(n\alpha)\Bigr|\le C_1\Bigl(\frac N{\log N}
  +\frac{N(\log R)^{3/2}}{R^{1/2}}\Bigr).
\]
\end{theorem}

In this section $\chi$ is a primitive character modulo $q\ge3$.

\begin{lemma}\label{lem:polya}
Let $B\ge0$ be such that $\bigl|\sum_{m\le H}\bar\chi(m)e(m\beta)/m\bigr|\le B$ for every integer
$H\le q/2$ and every $\beta\in\R$. Then $\bigl|\sum_{n\le T}\chi(n)\bigr|\le\sqrt q\,(B+6)$ for every
$T\ge0$.
\end{lemma}

\begin{proof}
Let $\tau=\sum_{a\bmod q}\bar\chi(a)e(a/q)$. Since $\bar\chi$ is primitive,
$\sum_a\bar\chi(a)e(an/q)=\chi(n)\tau$ for every $n$, and summing the squares of these identities over
$n\bmod q$ gives $q\varphi(q)=|\tau|^2\varphi(q)$, so $|\tau|=\sqrt q$. Summing over $n\le T$,
\[
  \tau\sum_{n\le T}\chi(n)=\sum_{0<|a|\le q/2}\bar\chi(a)\,
  \frac{e(aT/q)-1}{1-e(-a/q)},
\]
where $a$ runs over the integers with $-q/2<a\le q/2$, $a\ne0$. For $0<|\theta|\le1/2$ put
$r(\theta)=(1-e(-\theta))^{-1}-(2\pi i\theta)^{-1}$; then $|1-e(-\theta)|=2|\sin\pi\theta|\ge4|\theta|$.
If $2\pi|\theta|\le1$, then $|e(-\theta)-1+2\pi i\theta|\le(2\pi\theta)^2$ and
$|r(\theta)|\le(2\pi\theta)^2/(4|\theta|\cdot2\pi|\theta|)<3$; otherwise
$|r(\theta)|\le1/(4|\theta|)+1/(2\pi|\theta|)<3$. The terms with $r(a/q)$ contribute at most
$6q$. Writing $a=\pm m$ with $1\le m\le q/2$, the remaining terms are $q/(2\pi i)$ times a
combination, with coefficients of modulus at most $1$, of four sums
$\sum_{m\le H}\bar\chi(m)e(m\beta)/m$ with $H\le q/2$, so they contribute at most $2qB/\pi$. Dividing
by $|\tau|$ gives the claim.
\end{proof}

\begin{lemma}\label{lem:smooth}
For $N\ge1$, $\sum_{p\le N}1/p\le\log\log N+10$. For $Y\ge2$ and $H\ge1$,
\[
  \sum_{\substack{n\le H\\P(n)\le Y}}\frac1n\le e^{11}\log Y.
\]
\end{lemma}

\begin{proof}
Legendre's formula gives $\log N!\ge\sum_{p\le N}\lfloor N/p\rfloor\log p$, and $\log N!\le N\log N$,
$\sum_{p\le N}\log p\le N\log4$; hence $\sum_{p\le N}\log p/p\le\log N+\log4$, and partial
summation gives the first bound. The sum in the second is at most
$\prod_{p\le Y}(1-1/p)^{-1}\le\exp\bigl(\sum_{p\le Y}\frac1{p-1}\bigr)$, and
$\sum_p\bigl(\frac1{p-1}-\frac1p\bigr)\le1$.
\end{proof}

\begin{lemma}\label{lem:major}
There is an absolute constant $C_2$ such that for $0<r<q$, $b\in\mathbb Z$, $\beta\in\R$, integers
$Y\ge\max(2,r)$ and $t\ge1$,
\[
  \Bigl|\sum_{\substack{n\le t\\P(n)>Y}}\chi(n)e\Bigl(n\Bigl(\frac br+\beta\Bigr)\Bigr)\Bigr|
  \le C_2\,\frac{r\log(qr+2)\,(2+2\pi|\beta|t)\,t}{Y^{1/20}\log Y}.
\]
\end{lemma}

\begin{proof}
Let $\xi$ be a nonprincipal character modulo $k$. Proposition~\ref{prop:riesz} with
$\delta=1/2000$ and $a=7/8-2\delta$, which applies by Theorem~\ref{thm:zerofree}, gives
$|S(\xi;y)|\le Ky^{7/8+6\delta}\log(k+2)$ for $y\ge1$. Since
$(x+h)S(\xi;x+h)-xS(\xi;x)=h\sum_{n\le x}\xi(n)\Lambda(n)+\sum_{x<n\le x+h}\xi(n)\Lambda(n)(x+h-n)$,
the choice $h=x^{47/50}$, $\log x\le100x^{1/100}$ and the removal of the prime powers give
$\bigl|\sum_{p\le x}\xi(p)\log p\bigr|\le C_3x^{19/20}\log(k+2)$, and partial summation gives, for
integers $2\le A\le B$ and $\gamma\in\R$,
\begin{equation}\label{eq:twistedprime}
  \Bigl|\sum_{A<p\le B}\xi(p)e(p\gamma)\Bigr|\le
  2C_3\,\frac{B^{19/20}\log(k+2)(2+2\pi|\gamma|B)}{\log A}.
\end{equation}
Every $n\le t$ with $P(n)>Y$ is uniquely $n=pm$ with $p=P(n)$ prime, $Y<p\le t/m$ and
$P(m)\le p$, and then $m\le t/Y$. For such $p$ we have $p>r$, so $p\nmid r$, and on the units
modulo $r$ the function $u\mapsto e(umb/r)$ equals $\sum_\psi c_\psi\psi(u)$, the sum over the
characters $\psi$ modulo $r$, with $|c_\psi|\le1$. The character $\chi\psi$ modulo $qr$ is
nonprincipal: otherwise $\chi(n)=\bar\psi(n)$ for every $n$ prime to $qr$, so $\chi$ would be
induced by a character modulo $\gcd(q,r)<q$. Applying~\eqref{eq:twistedprime} to $\xi=\chi\psi$
with $\gamma=m\beta$, $A=\max(Y,P(m)-1)$ and $B=\lfloor t/m\rfloor$ gives
\[
  \Bigl|\sum_{\substack{n\le t\\P(n)>Y}}\chi(n)e\Bigl(n\Bigl(\frac br+\beta\Bigr)\Bigr)\Bigr|
  \le2C_3r\,\frac{\log(qr+2)(2+2\pi|\beta|t)}{\log Y}\sum_{m\le t/Y}\Bigl(\frac tm\Bigr)^{19/20},
\]
and $\sum_{m\le M}m^{-19/20}\le20M^{1/20}$.
\end{proof}

\begin{proposition}\label{prop:hsum}
There is an absolute constant $C_4$ such that if $\log q\ge200$, then
$\bigl|\sum_{m\le H}\chi(m)e(m\beta)/m\bigr|\le C_4\log\log q$ for every integer $H\le q$ and every
$\beta\in\R$.
\end{proposition}

\begin{proof}
Let $L=\log q$, $M=\lceil L\rceil$, $R=M^4$ and $Y=M^{200}$. Since $q=e^L\ge L^6/720$ and
$M\le L+1\le2L$, we have $2\le R<q$ and $R\le Y$. Split the sum according to $P(m)\le Y$ or
$P(m)>Y$. By Lemma~\ref{lem:smooth} the first part is at most $e^{11}\log Y$. For the second
put $A(t)=\sum_{n\le t,\,P(n)>Y}\chi(n)e(n\beta)$ and $w(t)=1/t-1/(t+1)$. Partial summation and
$|A(H)|\le H$ give
\[
  \Bigl|\sum_{\substack{m\le H\\P(m)>Y}}\frac{\chi(m)e(m\beta)}m\Bigr|\le1+\sum_{1\le t<H}|A(t)|\,w(t).
\]
If $t\le Y$, then $A(t)=0$. Let $t>Y$ and $Q=\lfloor t/R\rfloor\ge1$. By Dirichlet's theorem there
are coprime $a,r$ with $1\le r\le Q$ and $|\beta-a/r|\le1/(r(Q+1))$; then $|\beta-a/r|\,t\le R$.
If $r\le R$, then Lemma~\ref{lem:major} with $b=a$ gives $|A(t)|\le C_2kt$, where
\[
  k=\frac{R\log(qR+2)(2+2\pi R)}{Y^{1/20}\log Y}.
\]
If $r>R$, then $R\le r\le t/R$ and $|\beta-a/r|\le r^{-2}$, so Theorem~\ref{thm:mv} bounds the
full sum $\sum_{n\le t}\chi(n)e(n\beta)$, and removing the $\Psi(t)=\#\{n\le t:P(n)\le Y\}$ smooth
terms gives
\[
  |A(t)|\le C_1\Bigl(\frac t{\log t}+t\rho\Bigr)+\Psi(t),\qquad
  \rho=\frac{(\log R)^{3/2}}{R^{1/2}}.
\]
Since $tw(t)\le1/t$ and $w(t)\le1/t^2$,
\[
  \sum_{1\le t<H}|A(t)|\,w(t)\le C_1\sum_{t\le H}\frac1{t\log t}
  +(C_1\rho+C_2k)\sum_{t\le H}\frac1t+\sum_{1\le t<H}\Psi(t)w(t).
\]
The first sum is at most $3+\log L$, the second at most $1+L$, and by partial summation the third
is at most $\sum_{n\le H,\,P(n)\le Y}1/n\le e^{11}\log Y$. Finally $\log M\ge1$ and
$(\log M)^2\le4M$ give $\rho(1+L)\le32(\log M)^2/M\le128$; $\log(qR+2)\le6M$,
$2+2\pi R\le8M^4$ and $Y^{1/20}=M^{10}$ give $k(1+L)\le1$; and $\log Y=200\log M\le400\log L$.
Collecting terms, the sum is at most
$1+131C_1+C_2+(800e^{11}+C_1)\log L\le(800e^{11}+1+132C_1+C_2)\log L$.
\end{proof}

\begin{proof}[Proof of Theorem~\ref{thm:pv}]
The character $\bar\chi$ is primitive modulo $q$. If $\log q\ge200$, then
Proposition~\ref{prop:hsum} for $\bar\chi$ and Lemma~\ref{lem:polya} give
\[
  \Bigl|\sum_{n\le x}\chi(n)\Bigr|\le\sqrt q\,(C_4\log\log q+6)\le(C_4+6)\sqrt q\log\log q.
\]
If $\log q<200$, then every sum $\sum_{m\le H}\bar\chi(m)e(m\beta)/m$ with $H\le q$ has modulus at
most $1+\log q<201$, and Lemma~\ref{lem:polya} gives
\[
  \Bigl|\sum_{n\le x}\chi(n)\Bigr|\le207\sqrt q\le\frac{207}{\log\log3}\sqrt q\log\log q.\qedhere
\]
\end{proof}

\section{Formal verification and computations}\label{sec:lean}

The Lean~4 development (Lean v4.34.1, Mathlib v4.34.1) proves
{\footnotesize
\begin{verbatim}
theorem least_prime_primitive_root_bound {η : ℝ} (hη : 0 < η) :
    ∃ C : ℝ, 0 < C ∧ ∀ p : ℕ, p.Prime → ∃ ℓ : ℕ, ℓ.Prime ∧
      orderOf (ℓ : ZMod p) = p - 1 ∧
      (ℓ : ℝ) ≤ C * (((p - 1).divisors.card : ℝ) ^ 2 *
        Real.log (p + 2)) ^ (8 + η)
theorem least_prime_primitive_root_le_rpow {ε : ℝ} (hε : 0 < ε) :
    ∃ C : ℝ, 0 < C ∧ ∀ p : ℕ, p.Prime → ∃ ℓ : ℕ, ℓ.Prime ∧
      orderOf (ℓ : ZMod p) = p - 1 ∧ (ℓ : ℝ) ≤ C * (p : ℝ) ^ ε
theorem erdos_985_of_large : ∃ p₀ : ℕ, ∀ p : ℕ, p.Prime →
    p₀ ≤ p → ∃ q : ℕ, q.Prime ∧ q < p ∧
      orderOf (q : ZMod p) = p - 1
\end{verbatim}}
which are Theorem~\ref{thm:main} and Corollaries~\ref{cor:vinogradov} and~\ref{cor:erdos}, and
{\footnotesize
\begin{verbatim}
theorem lindelof {w : ℝ} (hw : 0 < w) :
    ∃ C : ℝ, 0 < C ∧ ∀ (q : ℕ) [NeZero q] (χ : DirichletCharacter ℂ q),
      χ ≠ 1 → ∀ s : ℂ, 7 / 8 + w ≤ s.re →
        ‖DirichletCharacter.LFunction χ s‖ +
          ‖(DirichletCharacter.LFunction χ s)⁻¹‖ ≤
          C * ((q : ℝ) * (3 + |s.im|)) ^ w
theorem character_sum_bound {w : ℝ} (hw : 0 < w) :
    ∃ C : ℝ, 0 < C ∧ ∀ (q : ℕ) [NeZero q] (χ : DirichletCharacter ℂ q),
      χ ≠ 1 → ∀ x : ℝ, 1 ≤ x →
        ‖∑ n ∈ Finset.Icc 1 ⌊x⌋₊, χ n‖ ≤
          C * (q : ℝ) ^ w * x ^ (7 / 8 + w)
theorem character_sum_le_rpow {η θ : ℝ} (hη : 0 < η) (hθ : 7 / 8 < θ) :
    ∃ C : ℝ, 0 < C ∧ ∀ (q : ℕ) [NeZero q] (χ : DirichletCharacter ℂ q),
      χ ≠ 1 → ∀ x : ℝ, 1 ≤ x → (q : ℝ) ^ η ≤ x →
        ‖∑ n ∈ Finset.Icc 1 ⌊x⌋₊, χ n‖ ≤ C * x ^ θ
\end{verbatim}}
\noindent(namespace \texttt{Vinogradov.CharSum}), which are Theorems~\ref{thm:lindelof}
and~\ref{thm:charsum} and Corollary~\ref{cor:charsum} (there $x\ge1$ follows from $x\ge q^\eta$),
and
{\footnotesize
\begin{verbatim}
def MVBound : Prop :=
  ∃ C : ℝ, 0 < C ∧ ∀ (q : ℕ) [NeZero q] (χ : DirichletCharacter ℂ q)
    (α : ℝ) (a : ℚ) (N : ℕ) (R : ℝ), 2 ≤ R → R ≤ a.den →
      (a.den : ℝ) ≤ N / R → |α - a| ≤ 1 / (a.den : ℝ) ^ 2 →
      ‖∑ n ∈ Finset.Ioc 0 N, (χ n : ℂ) * e (n * α)‖ ≤
        C * (N / Real.log N +
          N * Real.log R ^ (3 / 2 : ℝ) / Real.sqrt R)
theorem polya_vinogradov_loglog (hMV : MVBound) :
    ∃ C : ℝ, 0 < C ∧ ∀ (q : ℕ) [NeZero q] (χ : DirichletCharacter ℂ q),
      χ.IsPrimitive → 3 ≤ q → ∀ T : ℕ,
        ‖∑ n ∈ Finset.Ioc 0 T, (χ n : ℂ)‖ ≤
          C * Real.sqrt q * Real.log (Real.log q)
\end{verbatim}}
\noindent(namespace \texttt{Vinogradov.PV}, with $e(x)=\exp(2\pi ix)$), which is
Theorem~\ref{thm:pv}. \texttt{MVBound} is Theorem~\ref{thm:mv} for $f=\chi$, with the reduced
fraction $a/r$ given as a rational number $a$ with denominator \texttt{a.den}; it is a hypothesis
of \texttt{polya\_vinogradov\_loglog}.
Theorem~\ref{thm:zerofree} enters as the theorem
\[
  \texttt{OAI.DirichletCharacter.LFunction\_ne\_zero\_of\_seven\_eighths\_lt\_re}
\]
of~\cite{OAI26}, stated for Mathlib's \texttt{DirichletCharacter.LFunction} and compiled from its
source.
The Lean proofs of Lemmas~\ref{lem:logderiv} and~\ref{lem:lindelof} use the logarithmic
control lemma of~\cite{OAI26} in place of the Borel--Carath\'eodory and three circles arguments,
and the proof of
Lemma~\ref{lem:principal} uses Mathlib's bound $\psi(y)\ge(y-1)\log2-\log(y+2)$ in place
of~\cite{Nai82}. Every theorem depends only on the axioms \texttt{propext},
\texttt{Classical.choice} and \texttt{Quot.sound}.

Programs in Python, C, C++ and Julia compute $g^*(p)$ for all primes $3\le p\le10^8$
(Remark~\ref{rem:p0}), check Lemma~\ref{lem:indicator} numerically for $p\le200$, the inequality
of Lemma~\ref{lem:principal} on $10^4\le x\le2\cdot10^7$, Lemma~\ref{lem:arith} for
$N\le10^7$, the exponent identity in the proof of Theorem~\ref{thm:main} in exact arithmetic,
the two values in Remark~\ref{rem:carella}, the identity in the proof of
Lemma~\ref{lem:unsmooth} in exact arithmetic for six characters, the inequalities in the proof
of Proposition~\ref{prop:charsum} at $2\cdot10^5$ random points, the choices of parameters
in Section~\ref{sec:charsum} in exact arithmetic, $|\tau(\chi)|^2=q$ and
Lemma~\ref{lem:polya} for the quadratic characters modulo the primes $q\le3000$, the bound
$|r(\theta)|<3$ in its proof on a grid of $10^6$ points, and the inequalities between the
parameters in the proof of Proposition~\ref{prop:hsum}. For the quadratic characters modulo the
primes $10^3\le q\le2\cdot10^4$ the ratio $\max_x|\sum_{n\le x}\chi(n)|/(\sqrt q\log\log q)$
of Theorem~\ref{thm:pv} is at most $0.848$.

\section*{Data availability}
The Lean project, the programs and a script that runs them are archived at Zenodo,
doi:\href{https://doi.org/\zenodoID}{\zenodoID}.

\begin{thebibliography}{99}

\bibitem{BIMS17} V.~Bhargava, G.~Ivanyos, R.~Mittal and N.~Saxena, \emph{Irreducibility and
deterministic $r$-th root finding over finite fields}, arXiv:1702.00558, 2017.

\bibitem{Blo} T.~F. Bloom, \emph{Erd\H{o}s Problem \#985},
\url{https://www.erdosproblems.com/985}, accessed 14 October 2026.

\bibitem{Bur62} D.~A. Burgess, \emph{On character sums and primitive roots}, Proc. London
Math. Soc. (3) \textbf{12} (1962), 179--192.

\bibitem{Car15} N.~A. Carella, \emph{Least prime primitive roots}, Int. J. Math. Comput. Sci.
\textbf{10} (2015), no.~2, 185--194; arXiv:1709.01172.

\bibitem{Car18} N.~A. Carella, \emph{Primitive roots in short intervals}, arXiv:1806.01150v5,
2026.

\bibitem{Car25} N.~A. Carella, \emph{The least primitive roots mod $p$}, J. Math. Cryptol.
\textbf{19} (2025), article 20240017.

\bibitem{Dav00} H.~Davenport, \emph{Multiplicative Number Theory}, 3rd ed., revised by
H.~L. Montgomery, Grad. Texts in Math. \textbf{74}, Springer, New York, 2000.

\bibitem{Erd65} P.~Erd\H{o}s, \emph{Some recent advances and current problems in number theory},
in: Lectures on Modern Mathematics, Vol.~III, Wiley, New York, 1965, 196--244.

\bibitem{GS07} A.~Granville and K.~Soundararajan, \emph{Large character sums: pretentious
characters and the P\'olya--Vinogradov theorem}, J. Amer. Math. Soc. \textbf{20} (2007), 357--384.

\bibitem{GS18} A.~Granville and K.~Soundararajan, \emph{Large character sums: Burgess's
theorem and zeros of $L$-functions}, J. Eur. Math. Soc. \textbf{20} (2018), 1--14.

\bibitem{HW} G.~H. Hardy and E.~M. Wright, \emph{An Introduction to the Theory of Numbers},
6th ed., Oxford University Press, Oxford, 2008.

\bibitem{HB92} D.~R. Heath-Brown, \emph{Zero-free regions for Dirichlet $L$-functions, and the
least prime in an arithmetic progression}, Proc. London Math. Soc. (3) \textbf{64} (1992),
265--338.

\bibitem{Liu26} B.~Liu, \emph{Slightly improved zero-free half-planes for the
quasi-Riemann hypothesis}, arXiv:2610.12234, 2026.

\bibitem{Mar97} G.~Martin, \emph{The least prime primitive root and the shifted sieve}, Acta
Arith. \textbf{80} (1997), 277--288.

\bibitem{MV77} H.~L. Montgomery and R.~C. Vaughan, \emph{Exponential sums with multiplicative
coefficients}, Invent. Math. \textbf{43} (1977), 69--82.

\bibitem{Nai82} M.~Nair, \emph{On Chebyshev-type inequalities for primes}, Amer. Math. Monthly
\textbf{89} (1982), 126--129.

\bibitem{OAI26} OpenAI, \emph{The quasi-Riemann hypothesis: a zero-free half-plane
$\Re s>7/8$}, preprint, 30 September 2026, with a Lean~4 formalisation,
\url{https://github.com/openai/math}.

\bibitem{OAI26b} OpenAI, \emph{Primitive roots for every admissible integer base}, preprint,
4 October 2026, \url{https://github.com/openai/math}.

\bibitem{OAI26c} OpenAI, \emph{Deterministic polynomial factorization over prime fields},
preprint, 4 October 2026, \url{https://github.com/openai/math}.

\bibitem{Pal32} R.~E. A.~C. Paley, \emph{A theorem on characters}, J. London Math. Soc.
\textbf{7} (1932), 28--32.

\bibitem{Sho92} V.~Shoup, \emph{Searching for primitive roots in finite fields}, Math. Comp.
\textbf{58} (1992), 369--380.

\bibitem{Tit39} E.~C. Titchmarsh, \emph{The Theory of Functions}, 2nd ed., Oxford University
Press, Oxford, 1939.

\bibitem{Tit86} E.~C. Titchmarsh, \emph{The Theory of the Riemann Zeta-Function}, 2nd ed.,
revised by D.~R. Heath-Brown, Oxford University Press, Oxford, 1986.

\end{thebibliography}

\end{document}
